\documentclass[12pt,a4paper]{article}

% ---------- Мова та шрифти ----------
\usepackage{fontspec}
\usepackage{polyglossia}
\setmainlanguage{ukrainian}
\setotherlanguage{english}
\setmainfont{DejaVu Serif}
\setsansfont{DejaVu Sans}
\setmonofont{DejaVu Sans Mono}

% ---------- Компонування ----------
\usepackage[a4paper,margin=2.2cm]{geometry}
\usepackage{graphicx}
\usepackage{float}
\usepackage{caption}
\usepackage{xcolor}
\usepackage{hyperref}
\usepackage{enumitem}
\usepackage{titlesec}
\usepackage{longtable}
\usepackage{booktabs}
\usepackage{array}
\usepackage{amssymb}
\usepackage{fancyvrb}

\hypersetup{colorlinks=true, linkcolor=blue!50!black, urlcolor=blue!50!black,
    pdftitle={Порівняльний аналіз реалізацій --- Хмарне сховище}, pdfauthor={Проніна Марина}}

\titleformat{\section}{\normalfont\Large\bfseries}{\thesection.}{0.6em}{}
\titleformat{\subsection}{\normalfont\large\bfseries}{\thesubsection}{0.6em}{}
\setlength{\parindent}{0pt}
\setlength{\parskip}{0.6em}

\newcommand{\fig}[4]{%
    \begin{figure}[H]
        \centering
        \includegraphics[width=#3\linewidth,height=0.85\textheight,keepaspectratio]{#1}
        \caption{#2}
        \label{#4}
    \end{figure}
}

\begin{document}
\begin{titlepage}
\centering\large
\textbf{Київський національний університет імені Тараса Шевченка}\\
Факультет комп'ютерних наук та кібернетики
\vspace*{\stretch{1}}

\textbf{Звіт}\\
з дисципліни “Інформаційні Технології”\\
до Практичного завдання -- модуль 1\\
Варіант 1.2\\
“Етап 4. Порівняльний аналіз розробки”
\vspace*{\stretch{1}}

\begin{flushright}
Студентки 4 курсу\\
групи МІ-41\\
спеціальності 122 "Комп’ютерні науки"\\
Проніної Марини
\end{flushright}
\vspace*{\stretch{1.5}}

\textbf{Київ \the\year}
\end{titlepage}

\begin{abstract}
\noindent
У звіті порівнюються дві реалізації клієнта системи «Хмарне сховище» (тип 2, варіант 12): \textbf{десктоп-клієнт} на Python/PySide6 (етап 2) та \textbf{веб-клієнт} на Vue 3 (етап 3), що працюють з одним сервером FastAPI. Виділено спільні елементи й відмінності; порівняння охоплює архітектуру, функціональність, зручність, розгортання, продуктивність, тестованість і безпеку; наведено зведену таблицю, переваги й недоліки та обґрунтовано вибір найкращої технології.
\end{abstract}

\tableofcontents
\newpage


% =====================================================================
\section{Об'єкти порівняння}
% =====================================================================

Обидві реалізації використовують один і той самий REST API (\texttt{/\allowbreak{}api/\allowbreak{}auth/\allowbreak{}*}, \texttt{/\allowbreak{}api/\allowbreak{}files*}), одну базу та однакові правила варіанта 12 (перегляд \texttt{.xml}/\texttt{.png}, сортування за датою модифікації $\blacktriangle$/$\blacktriangledown$, фільтр «усі / .java / .png»). Отже, порівнюються саме клієнтські технології та способи доставки програми користувачу.

{\small
\begin{longtable}{@{}p{4cm}p{6cm}p{6cm}@{}}
\caption{Технологічний стек}\\
\toprule
\textbf{Аспект} & \textbf{Десктоп (етап 2)} & \textbf{Веб (етап 3)} \\
\midrule
\endhead
Мова & Python 3.11 & JavaScript (ES2022) \\ \midrule
GUI & PySide6 6.12 (Qt 6), MVVM & Vue 3.5 (Composition API), SFC \\ \midrule
Мережа & httpx (синхронно) & fetch (асинхронно) \\ \midrule
Синхронізація папки & \texttt{QFileSystemWatcher} + таймер, будь-яка ОС & File System Access API, лише Chromium \\ \midrule
Тести & pytest, 40 тестів & Vitest, 14 тестів \\ \midrule
Збірка / доставка & venv + pip (≈ 300 МБ з Qt) & \texttt{npm run build} $\to$ 87 КБ JS, роздається сервером \\ \midrule
\bottomrule
\end{longtable}}


% =====================================================================
\section{Спільні елементи та відмінності}
% =====================================================================

Спільні елементи обох реалізацій:

\begin{itemize}[nosep]
    \item один сервер і REST API (\texttt{/\allowbreak{}api/\allowbreak{}auth/\allowbreak{}*}, \texttt{/\allowbreak{}api/\allowbreak{}files*}), автентифікація JWT, ізольований простір користувача;
    \item однакові прецеденти: вхід/реєстрація, список файлів з атрибутами, показ/приховування стовпців, вивантаження, завантаження, видалення, drag-n-drop, перегляд \texttt{.xml}/\texttt{.png}, синхронізація папки;
    \item однакові правила варіанта 12 і однаковий тристоронній алгоритм синхронізації (\texttt{plan\_sync} / \texttt{planSync}) з розв'язанням конфліктів без втрати даних;
    \item розділення інтерфейсу та логіки (MVVM у десктопі, компоненти + модулі \texttt{lib/\allowbreak{}} у вебі) і модульні тести логіки.
\end{itemize}

{\small
\begin{longtable}{@{}p{4cm}p{6cm}p{6cm}@{}}
\caption{Відмінності реалізацій}\\
\toprule
\textbf{Аспект} & \textbf{Десктоп (етап 2)} & \textbf{Веб (етап 3)} \\
\midrule
\endhead
Середовище виконання & окремий процес Python + Qt на ПК & браузер, без встановлення \\ \midrule
Доставка й оновлення & встановлення на кожен ПК & роздається сервером, оновлюється централізовано \\ \midrule
Доступ до локальної папки & повний, події ОС (\texttt{QFileSystemWatcher}) & після вибору папки, лише Chromium, опитування \\ \midrule
Мережеві запити & синхронні (httpx) у GUI-потоці & асинхронні (fetch) \\ \midrule
Перегляд вмісту & бічна панель поруч із таблицею & модальне вікно \\ \midrule
Логіка варіанта & спільний з сервером модуль \texttt{core} & дубль мовою JavaScript (\texttt{variant.js}, \texttt{sync.js}) \\ \midrule
Мобільні пристрої & не підтримуються & адаптивна верстка \\ \midrule
\bottomrule
\end{longtable}}


% =====================================================================
\section{Порівняння за критеріями}
% =====================================================================

\subsection{Функціональність}

Обидва клієнти реалізують усі обов'язкові прецеденти й бонус drag-n-drop. Відмінність --- у синхронізації: десктоп-клієнт може стежити за папкою постійно та реагувати на зміни миттєво (події ОС), тоді як браузер надає доступ до папки лише після явного вибору користувачем і лише в Chrome/Edge, а відстежувати зміни може тільки опитуванням, поки вкладка відкрита.

\subsection{Зручність використання}

Веб-версія не потребує встановлення, адаптується до мобільних екранів і підтримує темну тему ОС; перегляд вмісту відкривається у модальному вікні. Десктоп-версія показує перегляд у бічній панелі поряд із таблицею, що зручніше при послідовному перегляді багатьох файлів, і інтегрується з провідником Windows (drag-n-drop, стандартні діалоги збереження).

\subsection{Розгортання та оновлення}

Веб-клієнт оновлюється для всіх користувачів одночасно: він входить в один Docker-образ із сервером, і кожен push у репозиторій автоматично перерозгортає його на Render. Десктоп-клієнт потребує встановлення Python-залежностей (або пакування, напр. PyInstaller) на кожному ПК і окремого оновлення.

\subsection{Продуктивність}

Для типових обсягів (десятки--сотні файлів) обидва клієнти працюють миттєво. Сортування й фільтрація виконуються локально в обох реалізаціях. Обчислення SHA-256 для синхронізації у Python (потоково, блоками 64 КБ) ефективніше за браузерне \texttt{crypto.subtle.digest}, яке вимагає завантажити весь файл у пам'ять. Десктоп-клієнт виконує мережеві запити синхронно в потоці GUI, що для великих файлів може на мить блокувати інтерфейс; веб-клієнт завдяки \texttt{async/\allowbreak{}await} цього недоліку не має.

\subsection{Тестованість і супровід}

Десктоп-клієнт ділить ядро (\texttt{cloudfolder.core}) із сервером --- правила варіанта існують у єдиному екземплярі. У веб-клієнті їх довелося продублювати мовою JavaScript; ризик розбіжностей знижено дзеркальними тестами з однаковими даними. В обох клієнтах логіку відокремлено від інтерфейсу (view-модель у десктопі, модулі \texttt{lib/\allowbreak{}} у вебі), тому вона тестується без GUI.

\subsection{Безпека}

В обох випадках автентифікація --- JWT (HS256, 12 год), паролі зберігаються як PBKDF2-SHA256 з сіллю, сервер ізолює простори користувачів і перевіряє імена файлів. Веб-клієнт зберігає токен у \texttt{sessionStorage} (видаляється із закриттям вкладки) і потенційно вразливий до XSS --- Vue екранує виведення за замовчуванням. Десктоп-клієнт тримає токен лише в пам'яті процесу.


% =====================================================================
\section{Зведена таблиця}
% =====================================================================

{\small
\begin{longtable}{@{}p{5cm}p{2cm}p{2cm}p{7cm}@{}}
\caption{Порівняння реалізацій (оцінка 1--5)}\\
\toprule
\textbf{Критерій} & \textbf{Десктоп} & \textbf{Веб} & \textbf{Коментар} \\
\midrule
\endhead
Повнота функцій варіанта & 5 & 5 & обидві реалізують усе, включно з drag-n-drop \\ \midrule
Синхронізація папки & 5 & 3 & веб --- лише Chromium і поки відкрита вкладка \\ \midrule
Встановлення / доступність & 2 & 5 & веб працює в браузері на будь-якому пристрої \\ \midrule
Оновлення & 2 & 5 & веб оновлюється централізовано \\ \midrule
Мобільні пристрої & 1 & 4 & адаптивна верстка \\ \midrule
Інтеграція з ОС & 5 & 3 & провідник, події файлової системи \\ \midrule
Чуйність інтерфейсу & 3 & 5 & синхронні запити в GUI-потоці десктопа \\ \midrule
Повторне використання коду & 5 & 3 & спільне ядро Python із сервером \\ \midrule
Тестованість & 4 & 4 & логіка тестується без GUI (pytest / Vitest) \\ \midrule
\textbf{Сума} & \textbf{32} & \textbf{37} &  \\ \midrule
\bottomrule
\end{longtable}}


% =====================================================================
\section{Переваги та недоліки}
% =====================================================================

\subsection{Десктоп-клієнт (PySide6)}

\begin{itemize}[nosep]
    \item \textbf{+} повноцінна фонова синхронізація з миттєвою реакцією на зміни в папці;
    \item \textbf{+} спільний з сервером код правил варіанта й синхронізації;
    \item \textbf{+} нативна інтеграція з ОС, перегляд у бічній панелі;
    \item \textbf{−} потрібне встановлення Python і Qt (великий обсяг), окреме оновлення кожного ПК;
    \item \textbf{−} не працює на мобільних пристроях; синхронні мережеві виклики в GUI-потоці.
\end{itemize}

\subsection{Веб-клієнт (Vue 3 SPA)}

\begin{itemize}[nosep]
    \item \textbf{+} нульове встановлення, доступ з будь-якого пристрою, централізоване оновлення;
    \item \textbf{+} асинхронний інтерфейс, адаптивна верстка, темна тема;
    \item \textbf{+} зручне налагодження інструментами розробника браузера;
    \item \textbf{−} синхронізація папки лише в Chromium і лише поки вкладка відкрита;
    \item \textbf{−} дублювання логіки варіанта мовою JavaScript.
\end{itemize}


% =====================================================================
\section{Обґрунтування найкращої технології}
% =====================================================================

За зведеною оцінкою (37 проти 32) найкращою для системи типу 2 є \textbf{веб-технологія: REST API на FastAPI + SPA на Vue 3}. Основне призначення системи --- доступ до особистої віддаленої папки з будь-якого місця, а саме тут веб має вирішальні переваги: відсутність встановлення, робота на будь-якому пристрої та централізоване оновлення. Єдина функція, в якій десктоп помітно сильніший, --- постійна фонова синхронізація; її доцільно залишити за легким десктоп-клієнтом як доповненням (так роблять і реальні сервіси на кшталт Google Drive). Обрана REST-архітектура дозволяє підтримувати обидва клієнти одним сервером без змін, тому рекомендована конфігурація --- веб-клієнт як основний інтерфейс і десктоп-клієнт як агент синхронізації.


% =====================================================================
\section{Висновки}
% =====================================================================

Обидві реалізації виконують усі вимоги варіанта 12 і проходять свої тести (40 тестів pytest і 14 тестів Vitest). Порівняння показало, що для системи віддаленої папки найкраще підходить веб-версія на REST + Vue 3, а десктоп-клієнт виправданий як спеціалізований інструмент синхронізації. Ключовим архітектурним рішенням, що уможливило дві реалізації без дублювання серверної логіки, стало відокремлення REST API від клієнтів.


\end{document}
